\section{Orbit quantization under stochastic compression}\label{sec:packet}
We first state the probabilistic argument in a real finite-dimensional
Hilbert space $H$. A complex unitary representation is covered by taking
its real inner product. Let $\mathbb S(H)$ be the unit sphere and let
$\mu$ be a probability measure on a compact group acting orthogonally
on $H$, with action $U_g$.

\begin{definition}[Orbit distortion]\label{def:distortion}
For $k\ge1$ set
\begin{equation}\label{eq:distortion}
 \Gamma_k(\mu)=\inf_{v,u_1,\ldots,u_k\in\mathbb S(H)}
  \int\min_{i\le k}\bigl(1-\langle U_gv,u_i\rangle\bigr)\,d\mu(g).
\end{equation}
It depends only on the representation and the law of a command packet,
not on a proposed stochastic realization. It lies in $[0,1]$.
\end{definition}
The upper bound one follows by using the direction of
$\int U_gv\,d\mu(g)$ as one of the centers, with any center if that
integral vanishes. The infimum is attained by compactness and continuity.
An invariant direction may make the coefficient zero; positivity is a
geometric property to be established, not an automatic conclusion for
all representations.

\begin{lemma}[One-packet contraction]\label{lem:packet}
Let $Y\in H$ be integrable, let $S$ be a finite register, and let $G$
be independent of $(Y,S)$ with law $\mu$. Suppose the new register
$S'$ has at most $k$ positive-probability values and its conditional
law, given $(Y,S,G)$, depends only on $(S,G)$. Then
\begin{equation}\label{eq:packet}
 \E\norm{\E[U_GY\mid S']}
 \le(1-\Gamma_k(\mu))\E\norm{\E[Y\mid S]}.
\end{equation}
\end{lemma}
\begin{proof}
Write $z_s=\E[Y\mid S=s]$, $p_s=\Pp(S=s)$ and
$T_g(s,r)=\Pp(S'=r\mid S=s,G=g)$. Conditional independence gives
\[
 p'_r z'_r=\sum_s p_s\int T_g(s,r)U_gz_s\,d\mu(g),
 \qquad z'_r=\E[U_GY\mid S'=r].
\]
Choose $u_r=z'_r/\norm{z'_r}$ when the latter is nonzero and any unit
vector otherwise. Thus
\begin{align*}
 \sum_r p'_r\norm{z'_r}
 &=\sum_s p_s\int\sum_rT_g(s,r)\langle U_gz_s,u_r\rangle\,d\mu(g)\\
 &\le\sum_{s:z_s\ne0}p_s\norm{z_s}
       \int\max_r\langle U_g(z_s/\norm{z_s}),u_r\rangle\,d\mu(g)\\
 &\le (1-\Gamma_k(\mu))\sum_s p_s\norm{z_s}.
\end{align*}
Zero centroids cause no loss of an identity: their inner product with
any chosen direction is zero. The maximization only relaxes the actual
possibly randomized assignment. In particular, neither the old register
nor the computation inside the packet is required to have $k$ states.
\end{proof}

\begin{theorem}[Representation-valued packet budget]\label{thm:group-budget}
Let $Y_0$ be a fixed unit vector and let successive command packets have
independent products $G_i$ with laws $\mu_i$. Between packets one may
insert fixed commands. A finite stochastic register processes the
commands in their prescribed order, with arbitrary cut-dependent rows.
Let $k_i$ bound its positive-probability states at the end of packet $i$.
If a terminal function $d(S_N)$ satisfies $\norm d\le B$ and
\[
             |\E\langle Y_N,d(S_N)\rangle|\ge B\kappa>0,
\]
then
\begin{equation}\label{eq:product-budget}
 \kappa\le\prod_i(1-\Gamma_{k_i}(\mu_i)),
 \qquad \sum_i\Gamma_{k_i}(\mu_i)\le\log(1/\kappa).
\end{equation}
Here $Y$ evolves by the same ordered representation products as the
commands. The statement also applies to a complex scalar correlation
with absolute value, using $q_N\ge\kappa$ as its terminal hypothesis.
\end{theorem}
\begin{proof}
At each boundary put $q=\E\norm{\E[Y\mid S]}$. A whole packet has
one stochastic input/output kernel conditional on its starting state
and its command word. Its internal rows and all internal hidden states
are integrated into that kernel. The fresh packet is independent of the
past, so Lemma~\ref{lem:packet} applies. A fixed command can only decrease
$q$, by conditional Jensen and preservation of norm by its action.
Initially $q=1$. At the terminal cut,
$|\E\langle Y_N,d(S_N)\rangle|\le Bq_N$.
Multiplication proves the first assertion; $1-u\le e^{-u}$ proves
the second. If a coefficient is one, positive terminal correlation is
impossible, consistent with the first assertion.
\end{proof}

The products $G_i$ need not commute. The packet endpoints need not be
consecutive cuts. Most importantly, the commands \emph{inside} a packet
need not be independent. Only its independence from the earlier process
is used in the conditional-mean identity. This is what permits a stronger
converse than a fair one-command test.

\begin{proposition}[A small-ball criterion]\label{prop:small-ball}
For $r>0$ define
\[
 A_\mu(r)=\sup_{v,u\in\mathbb S(H)}
              \mu\{g:\norm{U_gv-u}<r\}.
\]
Then
\begin{equation}\label{eq:smallball}
 \Gamma_k(\mu)\ge\frac{r^2}{2}\bigl(1-kA_\mu(r)\bigr)_+.
\end{equation}
If for every $v$ the packet is uniform on $m$ orbit points separated by
Euclidean distance at least $\sigma$, then
$\Gamma_k(\mu)\ge(1-k/m)_+\sigma^2/8$.
\end{proposition}
\begin{proof}
Outside the union of $k$ radius-$r$ balls the squared distance to each
center is at least $r^2$. Use
$1-\langle v,u\rangle=\norm{v-u}^2/2$ and the union bound.
For separated points, each open ball of radius $\sigma/2$ contains at
most one, so $A_\mu(\sigma/2)\le1/m$.
\end{proof}
